\documentclass[11pt]{article}
\usepackage{amsmath,amssymb}
\usepackage[margin=1in]{geometry}
\usepackage{hyperref}

\title{ORION Paper IV: Non-Escalating Scientific Authority under Content-Bound Evidence and Protected Evaluation}
\author{Working framework draft}
\date{August 2026}

\begin{document}
\maketitle

\begin{abstract}
Citation correctness and factual support do not by themselves establish scientific authority. ORION Paper IV studies a non-escalating authority pipeline in which retrieved or generated content remains proposal-level; evidence references bind exact content and provenance; attribution and semantic support are evaluated separately; candidate checking mechanisms must satisfy lineage and hostile-discrimination requirements; evaluator and holdout identity are protected; and missing evidence, compromised independence or contamination produces \texttt{CANNOT\_CHECK/BLOCK} rather than a confidence value that can be averaged into authority. Local authority-laundering falsifiers exercise content substitution, provenance identity, weak checkers, same-lane verification and checker chronology through the implemented kernel. This establishes mechanism semantics only. Whether the full pipeline materially lowers false scientific-authority promotion compared with strong source-aware verification systems remains externally \texttt{CANNOT\_CHECK}.
\end{abstract}

\noindent\textbf{Keywords:} scientific verification; provenance; attribution; evaluator integrity; abstention; AI assurance.

\section{Problem and claim boundary}
Paper IV does not claim citation-aware generation, claim verification, provenance tracking, source-aware factuality, claim-level auditability, benchmark contamination detection or evaluator locking as new ideas. Recent systems already address these components directly. The residual question is whether their combination can be turned into an explicit scientific-authority transition whose default under unresolved prerequisites is non-escalation.

\section{Evidence identity and source ownership}
An evidence record must resolve through host-owned content/provenance identity rather than trusting a reference string supplied by the candidate. A claim can be factually supported somewhere while still being assigned to the wrong source. ProvenanceGuard makes this cross-source conflation problem explicit and provides a strong source-aware verification baseline \cite{provenanceguard2026}.

Paper IV therefore separates at least: claim correctness, exact evidence identity, source ownership/attribution, semantic support or contradiction, and content provenance. A successful value on one axis cannot silently compensate for failure on another.

\section{Checker admissibility and protected evaluation}
A checker is not authoritative merely because it returns a score or a non-empty explanation. Its evidence, authoring/evaluation lineage, chronology and ability to discriminate hostile negatives matter. Claim-level auditability work motivates persistent semantic provenance \cite{aar2026}; ProvenAI further separates citation fidelity from behavioral influence \cite{provenai2026}. ORION treats these as proposal/evidence coordinates rather than final promotion authority.

Evaluation itself is an attack surface. RewardHackingAgents directly benchmarks evaluator tampering and held-out leakage \cite{rewardhacking2026}, while search-time contamination work shows that browsing agents may retrieve public benchmark answers during inference \cite{stc2026}. Evaluator/holdout identity and access telemetry must therefore be frozen prospectively.

\section{Non-escalating authority}
Let $A$ denote the authority state of a claim. Paper IV does not define promotion as a weighted confidence sum. Promotion requires registered prerequisites over content, attribution, support, checker admissibility, evaluator integrity and contamination status. If a required coordinate is unresolved or incompatible, the transition is blocked or returns \texttt{CANNOT\_CHECK}. This fail-closed behavior is the central candidate delta.

\section{External hostile evaluation}
The external benchmark includes both clean positive cases and attacks: correct fact/wrong source, content substitution, semantically close source conflation, cited-but-non-influential evidence, pooled support with wrong ownership, weak/same-lane/post-hoc checkers, search-time benchmark leakage, evaluator modification, holdout access and genuinely insufficient evidence.

The outcome-blind design is registered as \texttt{P4.protected-authority.v1}. The protocol freezes the attack taxonomy, source-aware baselines, authority-governance ablations, resource/exclusion rules, statistical analysis, access policy and result-figure/table definitions. The accompanying threat model and custody policy make the candidate, protected evaluator/holdout and host-verifier roles explicit. The primary design target is an absolute five-percentage-point reduction in false authority promotion relative to the strongest matched source-aware baseline while clean-case authority coverage remains non-inferior within five percentage points. These are prospective margins, not results.

Baselines include citation presence, pooled-evidence NLI, AttributionBench-style attribution evaluation \cite{attributionbench2024}, ProvenanceGuard-style source-aware verification, iterative retrieve-or-verify and current auditability systems. CLAIM-BENCH \cite{claimbench2025} is relevant for scientific claim-evidence comprehension. Candidate-caused failures/blocks remain in the denominator. Exact attack labels, evaluator/holdout hash, subject revision, source snapshots and baseline configs remain unbound until the independent host establishes \texttt{EXECUTION\_FROZEN}. Search/file/patch/evaluator access telemetry is retained.

\section{Results status}
The local hostile suite verifies current implementation semantics on exact cases. External false-promotion reduction, source-aware superiority and safety/coverage trade-offs remain \texttt{CANNOT\_CHECK}. Issue \#59 owns the protected hostile benchmark programme; this manuscript must not state external authority results until the final V1 execution bindings and independently held campaign exist.

\section{Threat model, limitations and reproducibility}
No finite hostile battery proves universal evaluator security. Human support/source judgments can be ambiguous. Web-search benchmarks may be contaminated even when local files are protected. The paper must report attack coverage, evaluator custody, access logs and correct abstentions, and must preserve all false positives/false negatives. The final artifact includes public portions of the attack manifest, content/provenance digests, baseline configs, per-claim verdicts, protected-access audit summaries and scripts regenerating every result figure/table. The frozen threat model, custody policy and normalized result schema are stored under \texttt{protocol/} and \texttt{research/paper-programme-v1/protocols/}.

\section{Conclusion}
Paper IV asks whether scientific authority should be an explicit protected state transition rather than an inference from fluent support scores. The mechanism is implemented; the external scientific claim remains open.

\bibliographystyle{plain}
\bibliography{bibliography}
\end{document}
